\section{编年事件}

\begin{event}{约公元前1046年}{牧野之战与商亡}{年代有争议}{牧野}
\term{这一年发生了什么}：周武王率联盟军在牧野击败商军，商纣王败亡。具体绝对年代在不同断代方案中存在差异，但周取代商的政治转折没有疑义。

\term{事情为什么会走到这里}：商末王权与诸方国关系紧张；周在渭水流域积累军事实力，并通过联盟取得东进条件。

\term{后来改变了什么}：周必须解决一个比战争更难的问题：如何统治旧商核心区及其众多附属政治体。
\end{event}

\begin{event}{约公元前1040年前后}{三监之乱与周公东征}{年代有争议}{东方旧商地区}
\term{经过}：武王死后，管叔、蔡叔等与武庚及东方势力发生反周叛乱；周公东征平定，并重组东方封建布局。

\term{前因}：新王朝继承、幼主摄政、旧商势力与周宗室内部权力关系同时不稳。

\term{后果}：周公摄政的合法性叙事、成周洛邑的建设以及对殷遗民的安置，成为西周早期秩序的关键支柱。
\end{event}

\begin{event}{公元前977年}{周昭王南征失利}{较可信}{汉水流域}
\term{经过}：周昭王南征，传世文献记其“丧六师于汉”，王本人亦未能返回。

\term{前因}：周王室试图向南方扩展资源与政治影响，南方地理、交通与地方政治条件远超西周核心区的控制能力。

\term{后果}：王室军事损失削弱了其对诸侯的威望；南方扩张的高成本成为后续王权资源约束的一部分。
\end{event}

\begin{event}{公元前841年}{国人暴动与共和行政}{可靠纪年}{镐京}
\term{经过}：厉王出奔，周、召二公共同执政。此年通常被视为中国传统连续纪年的确切起点。

\term{前因}：资源政策、政治排斥与王室—贵族关系失衡的累积。

\term{直接后果}：王位继承被暂时中断，行政由贵族集团共同维持。
\end{event}

\begin{event}{公元前771年}{犬戎攻破镐京，幽王死}{可靠纪年}{镐京}
\term{经过}：申侯与犬戎等攻入镐京，幽王死；次年平王东迁洛邑。

\term{前因}：幽王废申后与太子宜臼，改立褒姒之子伯服，触发申侯等政治联盟反弹；西部边防与王畿防御同时失效。

\term{后果}：西周结束。周王室迁往洛邑后仍保有名义共主地位，却失去西部王畿与直接军事实力，春秋时代的权力结构由此展开。
\end{event}

\begin{causalchain}[西周至东周]
分封网络曾让周以较低行政成本扩张；当王室继承危机和边防失效发生时，同一网络也让强大封国拥有独立行动空间。东迁不是单纯的地理移动，而是王室从资源中心退到礼仪中心的制度性降级。
\end{causalchain}
